% !TeX spellcheck = it_IT

\documentclass[aspectratio=169, handout]{beamer}

\usepackage[utf8]{inputenc}
\usepackage[italian]{babel}
\usepackage{hyperref}
\usetheme[titleformat=smallcaps, numbering=fraction, background=light, progressbar=frametitle]{metropolis}

\usepackage{mathtools}
\DeclarePairedDelimiter{\ceil}{\lceil}{\rceil}

\title{Esercitazioni di Informatica B}
\subtitle{Struct in MATLAB}

\author{Stefano Cereda\\
stefano.cereda@polimi.it
}
\date{26/11/2019}
\institute[PoliMi]{Politecnico Milano}
\logo{\includegraphics[width=10mm]{../logopolimi}}

\setbeamercovered{invisible}

\makeindex

\begin{document}
\begin{frame}
    \maketitle
\end{frame}

\begin{frame}{Array logici}
    Utilizzare gli array logici per risolvere i seguenti esercizi della scorsa esercitazione:
    \begin{itemize}
        \item Confronti --- slide 4
        \item Confronti con matrici --- slide 5
        \item Soglia intorno --- slide 10
        \item Caffè --- slide 11
        \item Matrice dominante --- slide 13
    \end{itemize}
\end{frame}

\begin{frame}{Paghe}
    Data la seguente tabella:

    \begin{tabular}{|c|c|c|c|c|c|}
        \hline
        & op1 & op2 & op3 & op4 & op5 \\
        \hline
        Paga oraria & 5 & 5.5 & 6.5 & 5 & 6.25 \\
        \hline
        Ore settimanali & 40 & 43 & 37 & 50 & 45 \\
        \hline
        Pezzi prodotti & 1000 & 1100 & 1000 & 1200 & 1100 \\
        \hline
    \end{tabular}

    Scrivere uno script Matlab che organizzi la tabella in un array di dati strutturati e risponda alle seguenti domande:
    \begin{enumerate}
        \item Quanto guadagna ogni operaio?
        \item Qual è il salario totale di tutti gli impiegati?
        \item Quanti pezzi vengono prodotti?
        \item Qual è il costo medio di un pezzo?
        \item Quante ore occorrono in media per un pezzo?
        \item Qual è l’operaio più efficiente?
    \end{enumerate}
\end{frame}

\begin{frame}[allowframebreaks]{Archivio mammiferi}

    Si vogliono rappresentare informazioni relative ai mammiferi ospitati in un
    parco protetto.

    In particolare, per ogni mammifero si rappresentano i seguenti dati:\\
    NOME TIPO, TIPO ALIMENTAZIONE, PESO MEDIO, insieme di ESEMPLARI

    Dove ogni esemplare è caratterizzato da:\\
    CODICE ESEMPLARE, ANNO DI NASCITA, NUMERO CUCCIOLI

    \begin{enumerate}
        \item Si crei una variabile mammifero che contenga le informazioni del
            seguente mammifero:\\
            NOME TIPO: Elefante\\
            TIPO ALIMENTAZIONE: Erbivoro\\
            PESO MEDIO: 2500\\
            ESEMPLARE 1: (Codice: 3, Anno: 1985, Cuccioli: 3)\\
            ESEMPLARE 2: (Codice: 5, Anno: 1989, Cuccioli: 0)\\
            ESEMPLARE 3: (Codice: 8, Anno: 1982, Cuccioli: 0)\\

        \item Chieda all'utente se vuole inserire ulteriori elefanti
        \item Si scriva il codice che costruisce, a partire dalla variabile creata al punto
            precedente, l'archivio degli esemplari di elefante sterili sapendo che un
            elefante si dice sterile se è nato prima del 1990 e non ha concepito
            nessun cucciolo.
    \end{enumerate}

    Si stampi su schermo l'elenco dei codici degli elefanti sterili.

\end{frame}

\begin{frame}[allowframebreaks]{Treni in stazione}
    Risolvere in MATLAB l'esercizio sulla stazione ferroviaria presentato nelle slide di riepilogo sul C.

    Lo script controlla se la stazione non è bloccata dalle manovre di qualche
    treno e se vi è almeno un treno in coda per entrare in stazione, quindi blocca la stazione, estrae dalla coda il
    primo treno che vi era stato inserito, gli assegna lo stato inIngresso, e aggiorna la coda in modo tale
    che poi contenga solo i rimanenti treni in coda.

    \break
    Se per esempio, prima di questa operazione la coda è cosi costituita: \\
    stazioneMI.codaTreni(1) = FR2092, stazioneMI.codaTreni(2) = IR097,
    stazioneMI.codaTreni(3) = FA2673 \\
    (per semplicità abbiamo riportato solo il nome del treno, ma, ovviamente, codaTreni dovrebbe contenere l’intera
    struct corrispondente), dopo l’operazione si deve avere che\\
    stazioneMI.codaTreni(1)= IR097, stazioneMI.codaTreni(2) = FA2673.
\end{frame}

\end{document}
